\section{Reproducing every number}
\label{app:repro}

No number in this paper is typed by hand. Two scripts read the result files and
emit the tables, the figures, and a file of \verb|\newcommand| macros that the
abstract and body quote, so the prose cannot drift away from the tables.

\subsection{Classical baselines}

The three solvers are pure NumPy and run on any machine in about six minutes,
most of it in the two spectral reference solves per material:

\begin{verbatim}
python wave/numerical/run_classical_benchmark.py \
    --out results/classical/classical_benchmark.json
\end{verbatim}

The material profiles used by the solvers are transcribed into NumPy so that
dataset generation does not require the torch stack;
\verb|material_profiles.verify_against_materials()| asserts they agree with the
project's originals to $10^{-12}$ and should be run before trusting any number
produced here.

\paragraph{Ground truth.} Chebyshev collocation at $n = 512$, with $n = 384$
run alongside as a self-convergence check. The two differ by
$\RefSelfConvHomogeneous\%$ on \emph{Homogeneous},
$\RefSelfConvTwoLayer\%$ on \emph{TwoLayer} and
$\RefSelfConvMultiLayer\%$ on \emph{MultiLayer} --- in every case far below the
$\FDFloor\%$ floor the second-order schemes plateau at, which is what makes the
floor measurable rather than an artefact of the reference.

\paragraph{Why not a fine finite-difference reference.} The obvious choice
would have been \texttt{fd2} at a very high resolution, and it is what the
project used previously. It does not work: \texttt{fd2}'s error stops falling
at $\FDFloor\%$, so scoring against a fine \texttt{fd2} solve measures
agreement with that floor rather than with the solution. The floor is identical
to three significant figures across three materials whose interior physics
differs completely, which is what identifies it as the boundary condition.

\subsection{Coordinate networks}

Runs are enumerated by an integer id; \verb|--list-runs| prints the full
mapping. A single run is

\begin{verbatim}
python wave/run_experiment.py --run-id 12 --output-dir results/
\end{verbatim}

and the ablation arms add \verb|--no-gradnorm| (pins
$w_{\mathrm{pde}} = w_{\mathrm{bc}} = w_{\mathrm{ic}} = 1$),
\verb|--causal-chunks 1| (one time slab, so every time is weighted equally),
and \verb|--run-tag| so an ablation does not overwrite its baseline.

\paragraph{Hyperparameters.} Identical across every architecture and every arm,
because the point of the comparison is the training recipe rather than a tuned
configuration: $100 \times 100$ collocation grid, $15\,000$ Adam steps at
$10^{-3}$, up to $9\,000$ continuation steps while $w_{\min} < 0.9$, then
$2\,000$ L-BFGS steps. Causal weighting uses $M = 16$ slabs with the tolerance
annealed from $0.1$ to $1.0$; GradNorm refreshes every $100$ steps with
smoothing $\alpha = 0.9$. Model selection is on a stationary validation metric
(unweighted per-slab residual plus boundary and initial-condition terms), not
the weighted training loss, which moves as the weights move.

\paragraph{Architecture sizes.} One representative configuration per
architecture, chosen so that the baseline, the SOAP/RBA arms and the ablation
arms all use the same network: 3 hidden layers of width 64 for VanillaPINN,
FourierFeature ($n_{\mathrm{F}} = 128$, $\sigma = 3$) and PirateNet; 2 hidden
layers of width 16 for the Kolmogorov--Arnold family (KAN $G = 5$, $k = 3$;
WavKAN Mexican hat; ChebyshevKAN degree 5; FourierWavKAN
$n_{\mathrm{F}} = 32$).

\paragraph{Reporting.} Each run writes \verb|l2_errors.json| (errors after the
Adam and L-BFGS phases, plus the arm's configuration) and
\verb|causal_convergence.json| (the per-slab weight vector sampled through
training, which is what Figure~\ref{fig:causal} plots). Both are small; the
network checkpoints are not needed to regenerate any number.

\subsection{Rebuilding the paper}

\begin{verbatim}
python wave/server/fetch_results.py          # fetch GPU results
python paper/neurips2026_workshop/make_tables.py
python paper/neurips2026_workshop/make_figures.py
pdflatex main && bibtex main && pdflatex main && pdflatex main
\end{verbatim}

\section{Implementation notes that changed a result}
\label{app:impl}

Three implementation choices were arrived at by fixing a failure, and all would
silently corrupt the comparison if made differently.

\paragraph{The spectral solver needs the first-order form.} Collocating
$u_{tt} = \rho^{-1}\partial_x(E\,\partial_x u)$ directly is the obvious route
and is unstable at the absorbing boundary: advancing $u$ at the end nodes from
$u^{n-1}$ decouples them from the interior and the corner modes grow without
bound within a few hundred steps. Splitting into velocity--stress makes the
radiation condition an exact statement about the incoming Riemann invariant,
which can be imposed as a projection at every Runge--Kutta stage. The
difference is not accuracy but existence: the second-order form produced
\texttt{NaN}, not a worse number.

\paragraph{Resampling can hide spectral convergence.} Solvers choose their own
time step from stability and their own spatial grid, so results are resampled
onto a shared $201 \times 101$ evaluation grid. Doing that with linear
interpolation is second-order accurate and puts a floor under the measured
error at roughly $10^{-3}$ --- which made the spectral solver look
\emph{algebraically} convergent and worse than finite differences. The fix is
in two parts: round the step count up until the output times are hit exactly,
removing time interpolation from the error budget entirely, and interpolate off
the Chebyshev nodes barycentrically rather than linearly. After both, the
spectral solver converges as it should. Any comparison of a spectral method
against a low-order one on a shared output grid is exposed to this.

\paragraph{Continuation must carry the final iterate.} We run continuation in
$3\,000$-step blocks so training can stop once $w_{\min}$ reaches its target.
Reloading the best validation checkpoint between blocks is not equivalent to
continuing: frontier progress need not improve that scalar metric, so the next
block can replay almost the same trajectory from the same parameters. The
runner now passes each block's final model directly to the next block and
restores the best-loss checkpoint only once, before L-BFGS. Result files record
this handoff mode, and every causal result reported in the paper was regenerated
after the fix.

\paragraph{A note on the KAN implementation.} The spline KAN used here is a
direct Equinox implementation of a B-spline edge layer with a SiLU base branch
(grid size $G$, spline order $k$), rather than a wrapper around an external
package. It matches the \verb|pykan| formulation but does not reproduce that
package's adaptive grid refinement, which would change the parameter shapes
mid-run --- neither the causal schedule nor L-BFGS tolerates that. The basis is
verified to form a partition of unity and to admit the second derivatives the
residual needs.
